\section{A core calculus, and what it buys}
\label{sec:formal}

\S\ref{sec:design} was operational, because that is how the system is built.
This section states the same design as a small calculus, so that the claims the mechanism rests on
become theorems rather than assertions: a freed slot is never dereferenced (\S\ref{sec:thm-safe}),
a compile-time tier switch is sound (\S\ref{sec:thm-tier}), removing collection work from an idle
function is observationally transparent (\S\ref{sec:thm-idle}), peak memory is a function of the live
set and not of the run length (\S\ref{sec:thm-bounds}), every synchronous collector step has a bound
that is independent of the heap (\S\ref{sec:thm-bounds}), total traversal work is linear in allocated
bytes (\S\ref{sec:thm-bounds}), and deferred freeing still makes progress
(\S\ref{sec:thm-bounds}).

Section~\ref{sec:formal-status} states exactly what is proved here, what is machine-checked in Lean,
what is checked by the compiler's own gates, and what is only tested.

\subsection{Store, frames, regions}
\label{sec:machine}

\emph{Frames} form a stack $F_0, \dots, F_n$ with $F_0$ outermost.
Each frame owns \emph{regions}, written $\rho^{\ell}$ where $\ell$ is the frame's distance from the
ground ($\rho^{0}$ belongs to $F_0$ and lives longest).
We use one direction throughout: \emph{smaller $\ell$ is longer-lived}.

A region is a sequence of \emph{slabs}; a slab is partitioned into \emph{slots}; a slot holds
either scalars, a struct, a \emph{view} $(p, \mathit{len})$, or a \emph{reference} $p$.
Every slot $\sigma$ has a type descriptor $\delta(\sigma)$ naming its reference words, and a
\emph{status} $\mathit{st}(\sigma) \in \{\mathsf{live}, \mathsf{free}\}$; the status of slots in a
released region is $\mathsf{free}$ and its memory is returned to the allocator.

Runtime state is a pair $(M, \mathit{sh})$ of a memory $M$ mapping slots to values and a shade set
$\mathit{sh}$ used by the incremental marker.
$\mathit{reach}_M(S)$ is the set of slots reachable from $S$ by following references and views
through the type descriptor of each slot; $\mathit{root}(F) = R(F)$ is the frame's \emph{root table},
a static list of the slots named by its reference-typed locals and parameters, emitted with the code
of $F$ and never computed at run time.
$\mathit{seeds}(\rho) = \bigcup_F R(F) \cup \mathit{foreign}(\rho)$ is the seed set of a region,
where $\mathit{foreign}(\rho)$ is the set of slots in $\rho$ that other regions register as
cross-region entry points.

\subsection{Static judgments}
\label{sec:judgments}

\paragraph{Placement.}
For an allocation site $s$ at depth $k$ the analysis yields a level
$\ell(s) = \max(0, k - \mathit{escapeDepth}(s))$ (eq.~\ref{eq:placement}) and the site allocates
directly in $\rho^{\ell(s)}$.
Nothing is moved later: there is no promotion copy at run time.

\paragraph{Depth rules.}
Every local, parameter and slot carries a \emph{depth} $d$ (the level of the region that holds it,
on the same scale; smaller is longer-lived):
\begin{enumerate}
\item[(D1)] a reference may be stored into a slot of depth $e$ only if the storage it names has
  depth $d \le e$;
\item[(D2)] if a value may escape $j$ frames, its site is placed at $k - j$, so at every escape site
  the reference is stored into a slot no shallower than the storage it names -- the store is one
  word write, and there is no copy of the value;
\item[(D3)] a view of container depth $d$ may only be used in a context whose depth $c$ satisfies
  $d \le c$.
\end{enumerate}
Violating (D1)--(D3) is a compile error, not a runtime check.
Two mutable references may alias and values are copied by default; no rule assumes uniqueness.

\paragraph{Tier judgments.}
$\mathit{acyclic}(U)$ holds when every reference write in $U$ either (i) stores a value freshly
allocated in the same frame, or a call result proved acyclic and fresh, (ii) stores \texttt{null},
(iii) stores a value from a strictly older level, or (iv) is a counted slot, and no other write can
drop a reference (this is the checker's \texttt{planFnAcyclic} condition, applied to every function
of the unit).
$\mathit{idle}(U)$ holds when every reference field is written at most once per object, no reference
local or parameter is reassigned in a way that drops a reference, and no operation drops a reference;
a call to a non-idle callee counts as a drop.

\subsection{Runtime transitions}
\label{sec:transitions}

The mechanisms of \S\ref{sec:contract} are transitions on $(M, \mathit{sh}, q, \mathcal{S})$,
where $q$ is the closing queue and $\mathcal{S}$ the free queue.
Write $\mathit{free}(\rho, S)$ for freeing the slots $S$ of $\rho$, and
$\mathit{live}(M) = \mathit{reach}_M(\bigcup_F R(F))$.

\begin{description}
\item[Allocate.] Take a free slot of the right size class (or bump a slab), set it $\mathsf{live}$,
  zero it, and if a mark round is in flight mark it black.
\item[Field store.] Write a value; if a mark round is in flight, shade the new value (insertion
  barrier) and, when replacing a reference, shade the old one (deletion barrier).
\item[Local release.] On block exit, the block's arena is rewound whole.
\item[Count zero.] A reference-counted slot whose count reaches zero is freed on the spot.
\item[Mark step.] Pop at most $b$ objects from the mark stack, scan each through its descriptor, and
  resume one chunked object scan if one is in progress; at most $\mathit{CHUNK}$ words of a resumed
  scan are visited before the cursor is saved again.
\item[Sweep step.] Visit at most $W_s$ words of a region's occupancy bitmap and free the slots of
  dead words.
\item[Trial delete.] On an overwrite that satisfies the four static conditions of \S\ref{sec:design},
  mark from $\mathit{seeds}(\rho) \setminus \{\sigma\}$ and free the slots reachable from the old
  value but unreachable from the remaining seeds.
\item[Close.] On region death: abort an in-flight round, drop the region's pins, set
  $\mathit{dead}(\rho)$, and append $\rho$ to $q$; this part is $O(1)$.
\item[Drain closing.] At an allocation point, pop regions from $q$, unlink their slots'
  cross-region edges, remove their registry entries, and push their slabs to $\mathcal{S}$;
  bounded by $\mathit{SLABS}$ slabs or $\mathit{D}_{	ext{cq}}$ microseconds per step.
\item[Drain free queue.] Return at most $\mathit{SLABS}$ queued slabs to the allocator, or stop at
  $\mathit{D}_{	ext{fq}}$ microseconds.
\item[Admit.] Before allocating $n$ bytes, if the next allocation would cross the trigger line, run
  the collector; under an explicit wall the trigger line satisfies $\mathit{thr} + R_{\max} \le M_{\max}$,
  so a request of size at most $R_{\max}$ always fits.
\end{description}

\subsection{Invariants and safety}
\label{sec:thm-safe}

\begin{description}
\item[(I1) Root completeness.] At every trigger point, $R(F)$ contains every live reference slot of
  $F$; for a region handed to the sweep, $\mathit{seeds}(\rho)$ contains every entry into $\rho$ from
  outside.
\item[(I2) Live references name live storage.] If $\sigma$ is live and its reference word holds a
  pointer to $\sigma'$, then $\sigma'$ is live.
\item[(I3) No live reference names freed storage.]
\item[(I4) Dead regions are unreachable.] After $\mathit{dead}(\rho)$ is set, no live slot and no
  seed names a slot of $\rho$.
\end{description}

\begin{assumption}[Escape and root soundness]
\label{as:sound}
(i) the placement and depth rules are enforced; (ii) at every trigger point the emitted root table
is complete in the sense of (I1); (iii) type descriptors list exactly the reference words of a layout;
(iv) C externs never write a reference word of a region slot.
\end{assumption}

\begin{lemma}[Local release is safe]
\label{lem:local}
Freeing a block's arena cannot free storage named by a live reference.
\end{lemma}
\begin{proof}
By (D1), a reference stored outside the block has depth at most the block's depth, so the object it
names has depth at most that too -- it is in this arena or a deeper one, all of which are released no
later. Hence no slot released here is named from an outer, still-live slot.
\end{proof}

\begin{lemma}[Sweep only frees unreachable slots]
\label{lem:sweep}
Under (I1), $\rho \setminus \mathit{reach}_M(\mathit{seeds}(\rho)) \subseteq
\bigcup_{\sigma \in \rho} \{\sigma \mid \text{unreachable}\}$.
\end{lemma}
\begin{proof}
$\mathit{live}(M)$ is reachable from the union of the frames' root tables, and by (I1) every entry into
$\rho$ from outside is in $\mathit{seeds}(\rho)$. A slot in $\rho$ that is reachable from a root is
therefore reachable from $\mathit{seeds}(\rho)$: take the last step of a path before it enters $\rho$
-- that step enters through an entry in the seed set. So the complement of the reachable set inside
$\rho$ is unreachable.
\end{proof}

\begin{lemma}[Trial deletion is conservative]
\label{lem:trial}
The set freed by a trial delete is a subset of the set a full sweep of $\rho$ with the same seeds
would free.
\end{lemma}
\begin{proof}
Trial deletion marks from the same seed set with $\sigma$ removed, so its reachable set is a superset
of the full sweep's reachable set; the freed set is the old value's reachable set minus a superset,
hence smaller. Removing $\sigma$ from the seeds is sound because $\sigma$ holds the old value at that
moment: including it would make the old structure its own justification.
\end{proof}

\begin{lemma}[Deferred region death keeps (I4)]
\label{lem:defer}
Setting $\mathit{dead}(\rho)$ at the owner frame's exit preserves (I2) and (I3), and makes
$\mathit{owns}(p) = 0$ for every $p \in \rho$.
\end{lemma}
\begin{proof}
At the exit of $F$, the only frames that may hold references into $\rho^{\ell}$ are $F$ and deeper
frames; the deeper ones have already exited, and their arenas are gone, so every remaining reference
into $\rho$ lives in $\rho$ itself (or in another region that is dying with the same frame).
Therefore no live slot names $\rho$, i.e. (I3) holds for it, and answering $0$ from
$\mathit{owns}$ cannot hide a live edge: any query that reaches a slot of $\rho$ started from dead
storage.
\end{proof}

\begin{theorem}[No dangling reference]
\label{thm:safe}
Under Assumption~\ref{as:sound}, no execution dereferences a freed slot.
\end{theorem}
\begin{proof}[Proof sketch]
Maintain (I1)--(I4) over the transitions of \S\ref{sec:transitions}.
(I1) is preserved by allocation (the allocating frame records the slot), by the store rule (a
reference stored into a slot of the current frame is a root, and a slot registered from another
region is a foreign seed), and by the fact that roots are only dropped together with the storage that
held them.
(I2) is preserved by allocate, store (the stored value is live by the induction hypothesis) and by
every rule that frees storage: local release (Lemma~\ref{lem:local}), sweep
(Lemma~\ref{lem:sweep}), trial delete (Lemma~\ref{lem:trial}).
(I3) is preserved by those lemmas plus deferred death (Lemma~\ref{lem:defer}), and by the drain
transition, which only unlinks cross-region edges of a region that is already dead.
Dereference is only possible through a reference stored in a live slot, and by (I2) that names a live
slot; a load through a stale pointer copied into a register earlier is excluded by (I3), since the
register contents came from a live slot at the time of the load.
\end{proof}

\subsection{Tier soundness}
\label{sec:thm-tier}

\begin{theorem}[Both tiers are sound, and the switch is a decision]
\label{thm:tier}
Let a unit $U$ be assigned a tier by \S\ref{sec:tiers}.
(a) If $\mathit{acyclic}(U)$, then in the event tier ($T1$) no slot is freed while reachable, and
every unreachable slot is freed within one reference-count update after it becomes unreachable.
(b) If $U$ is placed in the sweep tier ($T2$), then a sweep frees only unreachable slots, whether or
not reference counting is enabled.
\end{theorem}
\begin{proof}[Proof sketch]
(a) The conditions of $\mathit{acyclic}$ make the reference graph a DAG whose edges always point
from a younger (deeper) to an older (shallower) slot, or towards a fresh allocation. Counts
therefore decrease along a topological order: when the last external reference to an object is
dropped, every object reachable only through it has already had its own count reduced to the value
it would have if it were reachable, and a chain of zero counts is freed in one cascade.
No object can be freed while reachable because a reachable object has at least one holder on a live
path, and the cascade only starts from a slot whose count is zero.
The correctness of omitting the sweep is then immediate: the sweep would find every slot reachable
from the seeds, and every unreachable slot has already been freed by counting.
(b) The marker's seed set is the static root table plus the cross-region registration bits; by
Lemma~\ref{lem:sweep} the free set is a subset of the unreachable set. Counts play no role in the
seed construction, so the statement is independent of the \texttt{NOCOUNT} configuration; counting
on only adds the count-zero transition, which is sound by (I2) and can only make the reachable set
smaller.
\end{proof}

\begin{corollary}
\label{cor:tier-switch}
For an acyclic unit, switching the event tier off (forcing $T2$) cannot introduce a dangling
reference; it can only raise the peak memory, because more garbage waits for the sweep.
\end{corollary}

\subsection{Idle transparency}
\label{sec:thm-idle}

\begin{theorem}[Removing collection work from an idle function is transparent]
\label{thm:idle}
If $\mathit{idle}(f)$ holds, then emitting no collection work for $f$ leaves the observable trace of
the program unchanged; only timing changes.
\end{theorem}
\begin{proof}[Proof sketch]
The work emitted for a function is counting updates, budgeted collection steps, root materialization
for callees, and the store barrier.
The idle judgement removes the first three and keeps the barrier: a store an idle function does
perform still shades the edge it deletes, so the tri-colour invariant of the marker is maintained
exactly as in non-idle code.
Two facts are needed. First, a reference stored by an idle function either comes from a slot
reachable at the start of the cycle or is the result of an allocation in the current cycle (black by
construction): the idle conditions forbid reassigning reference locals and forbid writing a
reference field twice in one object, so the value's source is a root or an untouched field.
Second, no slot reachable at entry becomes unreachable during $f$, since $f$ drops no reference.
Hence no omitted step could have been the one that rescues a slot, and the region has the same live
set at exit.
Removing the steps therefore only removes work; induction on the statement structure, with a fixed
point over the call graph, makes the argument well-founded.
\end{proof}

\subsection{Bounds}
\label{sec:thm-bounds}

Let $L$ be the live set at a selection point, $G$ the trigger floor, $B$ the size of one slab and
$\kappa_M$ the measured excess over the live set per unit of retained garbage (the model of
\S\ref{sec:roofline}; $\kappa_M \in [0.772, 1.010]$).

\begin{theorem}[Bounded memory]
\label{thm:mem}
For a loop whose per-round allocation escapes into a bounded set of long-lived slots and whose other
allocations die with the round, the peak of a region satisfies
$M_{\text{peak}} \le L + G + O(B) = L(1 + \kappa_M \cdot r_{\max})$, independent of the number of
rounds, where $r_{\max} = G/L$ is the retention ratio the trigger line allows.
\end{theorem}
\begin{proof}[Proof sketch]
Round-local allocations are released by local release at block exit, so they contribute at most one
round's worth.
A promoted value either lands in a bounded slot set (contributing to $L$) or is overwritten; an
overwritten structure is freed by trial deletion as soon as the credit rule allows a traversal.
The credit rule requires unvisited allocation bytes to exceed $\max(2L, G)$, so between two
traversals the region grows by at most $G$ above its live set, and each denial postpones a free
without stacking: the next round's overwrite makes the older structure unreachable as well, and the
traversal eventually allowed frees the whole chain.
With an explicit wall, the trigger line additionally satisfies
$\mathit{thr} + R_{\max} \le M_{\max}$, so a request up to $R_{\max}$ bytes cannot overshoot the
wall by more than one request, and $R_{\max}$ is itself bounded by the largest single request the
program makes.
\end{proof}

\begin{theorem}[Bounded pause per step]
\label{thm:pause}
Every synchronous collector step performed on the mutator thread finishes in time bounded by
\begin{equation*}
\begin{aligned}
P_{\text{step}} &\le b \cdot \tau_t + \min(C_{\text{scan}}\cdot \mathit{CHUNK},\\
&\qquad D_{\text{seg}} + C_{\text{scan}}\cdot \mathit{CHK}) + W_s \cdot 64 \cdot \tau_s + c_0,
\end{aligned}
\end{equation*}
a constant fixed by compile-time budgets, independent of the heap size, of the number of live
objects, and of the size of the object being scanned; region destruction is $O(1)$ and its deferred
work is charged to the bounded drain steps.
\end{theorem}
\begin{proof}[Proof sketch]
Each step's work item is bounded by construction.
A mark step pops at most $b = \mathit{MARK\_BATCH\_MAX}$ objects and, if an object scan is in
progress, visits at most $\mathit{CHUNK}$ words or stops at the next multiple of
$\mathit{CHK}$ words after the segment deadline $D_{\text{seg}}$; a large object is therefore
scanned across several steps with a saved cursor rather than in one.
A sweep step visits at most $W_s$ bitmap words and frees at most $64 W_s$ slots; the finishing
credit for a slab is clamped by a constant so that the bound does not scale with the slab size.
A trial delete is bounded by the same mark budget because it consumes credit that the credit rule
only grants in proportion to allocated bytes.
Region death does constant work and appends to the closing queue; the queue is drained at
allocation points with at most $\mathit{SLABS}$ slabs or $D_{\text{cq}}$ microseconds per step, and
the free queue likewise.
Summing the bounded terms gives the stated constant; in the measured configuration
$b\cdot\tau_t \approx 4.5$\,\us, $C_{\text{scan}}\cdot\mathit{CHUNK} \approx 8$--$13$\,\us and
$W_s\cdot 64\cdot \tau_s \approx 0.23$\,\us, i.e. a few tens of microseconds per step, which is
what \S\ref{sec:eval} measures.
\end{proof}

\begin{theorem}[Amortized traversal work]
\label{thm:amortized}
Over an execution, the total traversal work charged to a region is
$O(\text{allocated bytes})$ for a region the checker proves acyclic (which additionally never
sweeps), and $O(\text{allocated bytes} + L)$ otherwise; the second bound constrains the total, not a
single step.
\end{theorem}
\begin{proof}[Proof sketch]
For the acyclic case the trigger is never armed with a non-empty mark phase, so the collector is a
no-op and the claim is immediate.
Otherwise each step's work is charged to allocation bytes admitted since the last visit by the credit
rule, and the credit only grows with allocation, so summing over the run gives the bound.
The bound deliberately says nothing about how the work is distributed over the steps; \S\ref{sec:thm-bounds}
supplies that, and \S\ref{sec:eval} measures the distribution.
\end{proof}

\begin{theorem}[Deferred freeing makes progress]
\label{thm:progress}
After a region $\rho$ is closed, the work needed to release it is $O(\text{slabs}(\rho))$, and each
allocation point performs at most $\mathit{SLABS}$ slabs of it, so if the program allocates at all,
$\rho$ is fully released after at most $\lceil \text{slabs}(\rho)/\mathit{SLABS} \rceil$ allocation
points; the closing queue holds at most $\mathit{CQ\_MAX}$ regions, and exceeding that forces an
immediate bounded drain.
\end{theorem}
\begin{proof}[Proof sketch]
Closing does constant work per region and enqueues it.
The drain processes a bounded number of slabs per step and never needs to revisit a slab: unlinking
the cross-region edges of a dead region is a finite walk over its occupied slots, and by
Lemma~\ref{lem:defer} none of those edges can be live.
Since $\mathit{dead}(\rho)$ makes $\mathit{owns}$ answer $0$ for all its slots, the region cannot
acquire new edges while it waits, so the remaining work is monotonically decreasing.
If the program stops allocating entirely before the queue is drained, the remaining slabs are
returned by the process exit path (the OS reclaims them); we state this as an explicit exception
because it is not a bound on the program's own memory.
\end{proof}

\subsection{Status of each claim}
\label{sec:formal-status}

\begin{table*}[t]
\centering
\small
\resizebox{\ifdim\width>\linewidth\linewidth\else\width\fi}{!}{\begin{tabular}{@{}llll@{}}
\toprule
Claim & Paper & Lean~4 & Gate \\
\midrule
Thm.~\ref{thm:safe} safety & proved (sketch + appendix) & \texttt{World.safe\_of\_invariants} & ASan/UBSan corpus \\
Thm.~\ref{thm:tier} tiers & proved (sketch) & \texttt{ExtCRegion}, \texttt{ExtCReclaim} & acyclic/forced red proofs \\
Thm.~\ref{thm:idle} idle & proved (sketch) & \texttt{ExtCIdleMax} (exact) & idle corpus ratchet \\
Thm.~\ref{thm:mem} memory & proved under Assumption~\ref{as:sound} & \texttt{ExtCGap}, \texttt{ExtCMemGap} & boundedness shapes \\
Thm.~\ref{thm:pause} pause & proved by budget counting & \texttt{ExtCPrecision} (footprint per slice) & work-budget gate \\
Thm.~\ref{thm:amortized} cost & proved (sketch) & -- & instruction-count gates \\
Thm.~\ref{thm:progress} deferred free & proved (sketch) & -- & closing-queue driver tests \\
\bottomrule
\end{tabular}}
\caption{What backs each claim. ``Proved'' means the argument in this section; ``Lean'' means the
machine-checked module that covers the corresponding core lemma; ``Gate'' means the automated check
that fails if the implementation stops matching. Blank cells are honest gaps, not claims.}
\label{tab:claims}
\end{table*}

\paragraph{What is assumed.}
Assumption~\ref{as:sound} is the boundary of the whole development.
(i) is the escape analysis, which is a large program that we test and gate but do not prove.
(ii) is what the compiler emits; it is checked structurally by gates that require every trigger point
to materialize the root set the analysis computed.
(iii) and (iv) are checked by construction (descriptors are emitted from the layout) and by the
extern convention (an extern receives references but never writes our slots).

\paragraph{What the Lean development covers.}
Thirteen modules, 2,555 lines, no \texttt{sorry}, checked by \texttt{proofs/verify.sh}: a region
calculus (\texttt{ExtCRegion}, \texttt{ExtCBlock}), the reclamation core
(\texttt{ExtCReclaim}, \texttt{ExtCRuntime}), the gap and memory laws used by
\S\ref{sec:roofline} (\texttt{ExtCGap}, \texttt{ExtCMemGap}), the idle maximum
(\texttt{ExtCIdleMax}), slice footprints (\texttt{ExtCPrecision}) and the pool and arena models.
The mapping is one-to-one for safety and idle; for the memory and pause bounds it covers the
underlying law rather than the whole implementation, and \tab\ref{tab:claims} says so.
Rebuilding the development from a clean tree compiles all thirteen modules, and the audit behind
\tab\ref{tab:claims} reproduces: no module contains \texttt{sorry}, and the only axioms any theorem
depends on are \texttt{propext} and \texttt{Quot.sound}.  Two practical notes for a reader repeating
this: the shipped script compiles modules in a fixed order and therefore needs the \texttt{.olean}
files of its dependencies to exist, so the modules are built in dependency order first
(\texttt{ExtCRegion}, then \texttt{ExtCPool}, \texttt{ExtCModel}, \texttt{ExtCPath}, and so on), and
its SMT cross-check step needs the \texttt{z3} Python package, which our environment does not have.

\paragraph{What is only tested.}
The distribution of pause times, the constants of the roofline model, and every number in
\S\ref{sec:eval} are measurements, not theorems.
Their reproduction is scripted, and the scripts are the artifact of \S\ref{sec:artifact}.
